\section{Stability of the subsolution threshold}
\label{sec:individual-density-stability}

We now study how the threshold $b_h(f)$ depends on $f$. By \eqref{eq:integrable-gauduchon-threshold}, $b_h^{1/n}$ is an infimum of functions that are continuous on $L^1$, and is therefore upper semicontinuous. It is not continuous in general: when $v_-(h)=0$, every density with positive threshold is an $L^1$ limit of densities with zero threshold. For density roots, \eqref{eq:integrable-gauduchon-root} expresses $\rho_h$ as a supremum of positive linear functionals. Hence $\rho_h$ is a norm on the space where it is finite, this space is complete for a suitable norm, and the threshold depends continuously on the root in that norm.

\begin{theorem}\label{thm:individual-threshold-semicontinuity}
Suppose $f_j,f$ are nonzero nonnegative $L^1$ densities and $f_j\to f$ in $L^1(\mathrm{d}V)$. Then
\[
 \limsup_{j\to\infty}b_h(f_j)\leq b_h(f).
\]
Moreover,
\begin{equation}\label{eq:individual-threshold-relaxation}
 b_h(f)=\lim_{\delta\downarrow0}\sup\{c_h(q):q\in L^\infty,\ q\geq0,\ q\not\equiv0,
                                      \ \|q-f\|_{L^1}<\delta\},
\end{equation}
and the suprema are finite for sufficiently small $\delta$.
\end{theorem}
\begin{proof}
For a Gauduchon form $\gamma$ on $X$, write $A_\gamma(f)=\int_X(f\,\mathrm{d}V/\gamma^n)^{1/n}\gamma^n$. Since $|s^{1/n}-t^{1/n}|\leq|s-t|^{1/n}$ for $s,t\geq0$, H\"older's inequality shows that $A_\gamma$ is continuous on $L^1$; it is positive at every nonzero nonnegative density. By \cref{thm:integrable-gauduchon-formula}, $b_h^{1/n}$ is the infimum over $\gamma$ of the continuous positive functions $f\mapsto\int_Xh\wedge\gamma^{n-1}/A_\gamma(f)$. It is therefore upper semicontinuous, and testing a single $\gamma$ bounds it above on an $L^1$ neighborhood of $f$.

For the second assertion, recall that $b_h(q)=c_h(q)$ for bounded $q$ (\cref{prop:individual-normalizer-lp}). By the first assertion, $\limsup_jc_h(q_j)\leq b_h(f)$ for every sequence of nonzero nonnegative bounded densities $q_j\to f$ in $L^1$. Hence the limit in \eqref{eq:individual-threshold-relaxation} is at most $b_h(f)$, and the local upper bound makes the suprema finite for small $\delta$. Conversely, the truncations $f_k$ converge to $f$ in $L^1$ and satisfy $c_h(f_k)\geq b_h(f)$, so every supremum in \eqref{eq:individual-threshold-relaxation} is at least $b_h(f)$.
\end{proof}

Since $b_h\geq0$, upper semicontinuity makes $b_h$ continuous at every density where it vanishes. If $v_-(h)=0$, these are its only points of continuity. Indeed, let $b_h(f)>0$ and $g=f/\int_Xf\,\mathrm{d}V$. By \cref{cor:lower-mass-lone-instability} and \cref{thm:individual-density-subsolution}, the densities $(1-t)g+tf_0$, where $f_0$ is the density of \cref{thm:lower-mass-obstructing-measure}, have zero threshold and converge to $g$ in $L^1$ as $t\downarrow0$; multiplying them by $\int_Xf\,\mathrm{d}V$ gives densities with zero threshold converging to $f$. These statements concern only the scalar $b_h$; they say nothing about capacity convergence of the corresponding potentials or about the limits of their Monge--Amp\`ere measures.

\subsection{Finite sums of densities}\label{subsec:finite-density-sums}

Since \eqref{eq:integrable-gauduchon-root} expresses $\rho_h(p)$ as a supremum of positive linear functionals evaluated at $|p|$, the root functional is homogeneous, monotone and subadditive. Subadditivity shows that the sum of two densities with positive threshold again has positive threshold. In the statements below, $f$ and $g$ are nonzero nonnegative $L^1$ densities, and $p$ and $q$ are real functions in $L^n$.

\Needspace{12\baselineskip}
\begin{proposition}\label{prop:individual-root-triangle}
The functional $\rho_h$ satisfies $\rho_h(tp)=|t|\rho_h(p)$ for real $t\ne0$,
$\rho_h(p)\leq\rho_h(q)$ whenever $|p|\leq|q|$, and
\[
 \rho_h(p+q)\leq\rho_h(p)+\rho_h(q).
\]
In particular, if $b_h(f),b_h(g)>0$, then
\[
 b_h(f+g)\geq
       \bigl(b_h(f)^{-1/n}+b_h(g)^{-1/n}\bigr)^{-n}>0.
\]
\end{proposition}
\begin{proof}
By \eqref{eq:integrable-gauduchon-root}, $\rho_h(p)=\sup_\gamma L_{h,\gamma}(|p|)$ is a supremum of positive linear functionals evaluated at $|p|$. This gives homogeneity and monotonicity, and, together with $|p+q|\leq|p|+|q|$, the triangle inequality. Finally, $(f+g)^{1/n}\leq f^{1/n}+g^{1/n}$ gives
\[
 b_h(f+g)^{-1/n}=\rho_h\bigl((f+g)^{1/n}\bigr)\leq\rho_h(f^{1/n})+\rho_h(g^{1/n})
 =b_h(f)^{-1/n}+b_h(g)^{-1/n}.
\]
\end{proof}

\begin{proposition}\label{prop:individual-density-banach}\index{Density space!completeness}
The real vector space
\[
 \mathcal D_h=\{p\in L^n(\mathrm{d}V):\rho_h(p)<+\infty\}
\]
is complete for the norm $\|p\|_{\mathcal D_h}=\|p\|_{L^n}+\rho_h(p)$. It is an order ideal in $L^n$: if $q\in\mathcal D_h$ and $|p|\leq|q|$, then $p\in\mathcal D_h$. It contains $L^\infty$ and is therefore dense in $L^n$ for the $L^n$ norm.

Moreover, $\mathcal D_h=L^n$ if and only if $v_-(h)>0$, and in that case the norm $\|\cdot\|_{\mathcal D_h}$ is equivalent to the $L^n$ norm. If $h$ is K\"ahler, then
\[
 \rho_h(p)=\left(\int_Xh^n\right)^{-1/n}\|p\|_{L^n}.
\]
\end{proposition}
\begin{proof}
By \cref{prop:individual-root-triangle}, $\mathcal D_h$ is a vector space and an order ideal, and $\|\cdot\|_{\mathcal D_h}$ is a norm on it. Bounded functions belong to $\mathcal D_h$, because $b_h(|p|^n)=c_h(|p|^n)>0$ for bounded nonzero $p$.

\Needspace{5\baselineskip}
We next show that $\rho_h$ is lower semicontinuous on $L^n$. If $p_j\to p$ in $L^n$, then $|p_j|^n\to|p|^n$ in $L^1$, and the upper semicontinuity of $b_h$ in \cref{thm:individual-threshold-semicontinuity} gives
\[
 \rho_h(p)\leq\liminf_j\rho_h(p_j).
\]
This includes the case $b_h(|p|^n)=0$, in which the right side must be infinite; for $p=0$ the inequality is trivial. Now let $(p_j)$ be a Cauchy sequence in $\mathcal D_h$. It converges in $L^n$ to some $p$, and lower semicontinuity gives $\rho_h(p)\leq\liminf_j\rho_h(p_j)<+\infty$, so $p\in\mathcal D_h$. Applied to $p_j-p_k$ as $k\to\infty$, lower semicontinuity also gives
\[
 \rho_h(p_j-p)\leq\liminf_{k\to\infty}\rho_h(p_j-p_k),
\]
which tends to zero as $j\to\infty$. Together with $L^n$ convergence, this proves completeness.

If $v_-(h)>0$, \cref{prop:individual-normalizer-lp} gives
$\rho_h(p)\leq v_-(h)^{-1/n}\|p\|_{L^n}$, so $\mathcal D_h=L^n$ with an equivalent norm. If $v_-(h)=0$, the density $f$ of \cref{thm:lower-mass-obstructing-measure} has $b_h(f)=0$ by \cref{thm:individual-density-subsolution}; thus $f^{1/n}$ belongs to $L^n$ but not to $\mathcal D_h$. The K\"ahler formula follows from the same proposition.
\end{proof}

Thus, when the lower mass vanishes, $\mathcal D_h$ is a proper dense subspace of $L^n$, complete for the stronger norm $\|\cdot\|_{\mathcal D_h}$; its nonnegative elements are the $n$th roots of the densities admitting a positive scaled subsolution. A finite sum of such densities again admits a positive scaled subsolution, but a convergent infinite sum need not: the smooth summands in \cref{thm:lower-mass-obstructing-measure} give an example. By \cref{prop:individual-threshold-backgrounds}, the space $\mathcal D_h$ does not depend on the Hermitian form $h$, and the norms $\|\cdot\|_{\mathcal D_h}$ for different Hermitian forms are equivalent. The closure of the bounded functions for this norm is studied in \cref{subsec:bounded-root-approximation}.

The restriction of $\rho_h$ to $\mathcal D_h$ is itself a norm: if $p\ne0$, then $\rho_h(p)\geq L_{h,\gamma}(|p|)>0$ for any Gauduchon form $\gamma$. Whether this weaker norm is complete depends on the upper mass.

\Needspace{15\baselineskip}
\begin{proposition}\label{prop:individual-root-completeness}\index{Density space!root norm}
The norm $\rho_h$ is complete on $\mathcal D_h$ if and only if $v_+(h)<+\infty$. When the upper mass is finite,
\[
 \|p\|_{L^n}\leq v_+(h)^{1/n}\rho_h(p)\qquad(p\in\mathcal D_h),
\]
so $\rho_h$ and $\|\cdot\|_{\mathcal D_h}$ are equivalent norms.
\end{proposition}
\begin{proof}
Suppose $v_+(h)<+\infty$. A continuous solution for a bounded density $q$ has mass $c_h(q)\int_Xq\,\mathrm{d}V\leq v_+(h)$. Applying this to the truncations of $|p|^n$ and letting the truncation level increase gives $b_h(|p|^n)\|p\|_{L^n}^n\leq v_+(h)$, which is the asserted inequality. Hence $\rho_h$ and $\|\cdot\|_{\mathcal D_h}$ are equivalent, and completeness follows from \cref{prop:individual-density-banach}.

Suppose instead that $v_+(h)=+\infty$. By \cref{thm:gaud-volume-mass-extrema,thm:gaud-volume-dual-formula}, the normalization scalars of smooth strictly positive probability densities are unbounded above. For each $j$, choose such a density $f_j$ with
\[
 c_h(f_j)\geq(j2^j)^n,
 \qquad\text{and put}\qquad p_j=jf_j^{1/n}.
\]
Then $\|p_j\|_{L^n}=j$ and $\rho_h(p_j)=jc_h(f_j)^{-1/n}\leq2^{-j}$, so the partial sums $P_N=\sum_{j\leq N}p_j$ form a Cauchy sequence for $\rho_h$.

Fix a Gauduchon form $\gamma$. Its coefficient $a_\gamma$ is bounded below by a positive constant, so \eqref{eq:integrable-gauduchon-root} gives $\rho_h(p)\geq c\|p\|_{L^1}$ with a constant $c>0$. Hence $\sum_jp_j$ converges in $L^1$ to a nonnegative function $P$. If $P_N$ converged for $\rho_h$ to an element of $\mathcal D_h$, that element would also be the $L^1$ limit, hence equal to $P$ almost everywhere. But $P\geq p_j$ gives $\|P\|_{L^n}\geq j$ for every $j$, so $P\notin L^n$. Therefore $\rho_h$ is not complete.
\end{proof}

The two mass properties play separate roles: positive lower mass makes every $L^n$ function admissible, whereas finite upper mass makes the root norm control the $L^n$ norm. The norm $\|\cdot\|_{\mathcal D_h}$ is complete without either assumption.

All four combinations occur. On a compact K\"ahler manifold both mass properties hold. The Hopf threefold has positive lower mass and infinite upper mass (\cref{cor:hopf-lower-mass}). The elliptic threefold of \cref{thm:vanishing-lower-mass} has zero lower mass, and finite upper mass for its balanced metric (\cref{prop:lower-mass-balanced-metric}). The product manifolds of \cref{cor:both-mass-properties-fail} have zero lower mass and infinite upper mass. Since neither $\mathcal D_h$ nor the completeness of $\rho_h$ depends on the choice of $h$ (\cref{prop:individual-threshold-backgrounds,prop:bmp-metric}), we obtain the following table.
\begin{center}
\begin{tabular}{@{}l@{\qquad}l@{\qquad}l@{}}
Manifold & Domain $\mathcal D_h$ & Is $\rho_h$ complete?\\
\hline
K\"ahler manifold & $L^n$ & Yes\\
Hopf threefold & $L^n$ & No\\
Elliptic threefold & $\mathcal D_h\subsetneq L^n$ & Yes\\
Product example & $\mathcal D_h\subsetneq L^n$ & No
\end{tabular}
\end{center}

\Needspace{14\baselineskip}
\subsection{Approximation by bounded density roots}\label{subsec:bounded-root-approximation}

Bounded functions are dense in $\mathcal D_h$ for the $L^n$ norm, but not obviously for the norm $\|\cdot\|_{\mathcal D_h}$. Let $\mathcal D_h^0$ denote the closure of $L^\infty$ in $\mathcal D_h$ for this norm. We show that a function belongs to $\mathcal D_h^0$ exactly when its tails are small uniformly over all Gauduchon tests. The basic tool is a subcritical estimate. For every $r>n$,
\begin{equation}\label{eq:root-subcritical-bound}
 \rho_h(p)\leq C_r\|p\|_{L^r}\qquad(p\in L^r),
\end{equation}
where the constant $C_r$ depends on $h$, the reference volume and $r$, but not on $p$. Indeed, for nonzero $p\in L^r$ the density $|p|^n$ belongs to $L^{r/n}$, so $b_h(|p|^n)=c_h(|p|^n)$ by \cref{prop:individual-normalizer-lp}, and the lower-density estimate of \cref{thm:lp-potential-estimate}, with exponent $r/n$, gives $c_h(|p|^n)\||p|^n\|_{L^{r/n}}\geq c_{r/n}>0$.

\begin{proposition}\label{prop:root-uniform-tails}\index{Uniform tails!root density}
For $p\in\mathcal D_h$, the following conditions are equivalent:
\begin{enumerate}
\item $p\in\mathcal D_h^0$;
\item $\rho_h(|p|\mathbf1_{\{|p|>K\}})\to0$ as $K\to\infty$;
\item whenever $|q_j|\leq|p|$ and $q_j\to0$ in measure, one has $\|q_j\|_{\mathcal D_h}\to0$.
\end{enumerate}
By \eqref{eq:integrable-gauduchon-root}, the second condition can be written as
\[
 \sup_{\gamma\text{ Gauduchon}}
       L_{h,\gamma}(|p|\mathbf1_{\{|p|>K\}})\longrightarrow0
 \qquad(K\to\infty).
\]
The space $\mathcal D_h^0$ is a closed order ideal in $\mathcal D_h$. It contains $L^r$ for every $r>n$, and smooth functions are dense in it. If $v_-(h)>0$, then $\mathcal D_h^0=\mathcal D_h=L^n$.
\end{proposition}
\begin{proof}
For $K>0$, set $p_K=\operatorname{sgn}(p)\min(|p|,K)$. Then $|p-p_K|=(|p|-K)_+$ and
\[
 (|p|-K)_+\leq |p|\mathbf1_{\{|p|>K\}},\qquad
 |p|\mathbf1_{\{|p|>2K\}}\leq2(|p|-K)_+.
\]
Since the $L^n$ norms of the tails tend to zero, these inequalities and the monotonicity of $\rho_h$ show that the second condition is equivalent to $p_K\to p$ in $\mathcal D_h$; in particular, it implies the first. Conversely, if $a$ is bounded and $K\geq\|a\|_{L^\infty}$, then $(|p|-K)_+\leq|p-a|$, so $\|p-p_K\|_{\mathcal D_h}\leq\|p-a\|_{\mathcal D_h}$. Approximating $p$ by bounded functions $a$ shows that the first condition implies $p_K\to p$, hence the second.

Assume the second condition, and let $q_j$ be as in the third. Fix $K$. On $\{|p|\leq K\}$, the functions $q_j$ are bounded by $K$ and tend to zero in measure, hence in $L^{2n}$; by \eqref{eq:root-subcritical-bound} with $r=2n$, the root functional of $q_j\mathbf1_{\{|p|\leq K\}}$ tends to zero. On the complementary set, $\rho_h(q_j\mathbf1_{\{|p|>K\}})\leq\rho_h(|p|\mathbf1_{\{|p|>K\}})$. Letting first $j$ and then $K$ tend to infinity gives $\rho_h(q_j)\to0$, and dominated convergence gives $\|q_j\|_{L^n}\to0$. Conversely, the third condition applied to $q_j=|p|\mathbf1_{\{|p|>j\}}$ gives the second.

Since the second condition passes from $p$ to every function dominated by $|p|$, the space $\mathcal D_h^0$ is an order ideal; it is closed by definition. For $p\in L^r$ with $r>n$, the truncations $p_K$ converge to $p$ in $L^r$, hence in $\mathcal D_h$ by \eqref{eq:root-subcritical-bound}. Bounded functions can be approximated in $L^{2n}$ by smooth functions, which then converge in $\mathcal D_h$ by the same estimate; hence smooth functions are dense in $\mathcal D_h^0$. Under positive lower mass, \cref{prop:individual-density-banach} shows that $\mathcal D_h=L^n$ with an equivalent norm, and bounded functions are dense in $L^n$; hence $\mathcal D_h^0=\mathcal D_h=L^n$.
\end{proof}

By \cref{prop:individual-threshold-backgrounds}, the space $\mathcal D_h^0$ does not depend on the choice of the Hermitian form $h$. Convergence in $\mathcal D_h$ controls the threshold: if $p_j\to p\ne0$ in $\mathcal D_h$, then
\[
 |\rho_h(p_j)-\rho_h(p)|\leq\rho_h(p_j-p)\longrightarrow0,
 \qquad b_h(|p_j|^n)\longrightarrow b_h(|p|^n)>0.
\]
For $p\in\mathcal D_h^0$, the uniform tail condition also gives a uniform modulus for exponential damping and error bounds for the scalars of the exponential equations; see \cref{prop:exponential-root-tail-modulus}. It does not by itself give capacity convergence of the solutions. We do not know whether $\mathcal D_h^0=\mathcal D_h$ when $v_-(h)=0$.

